% Teil II: durchgerechneter Weg, Teile a--c.

\subsection*{Gegeben und gesucht}
Kurze, in eigenen Worten zu beantwortende Verständnisfragen: (a) Inhalt
der drei Axiome, (b) Astronaut in der ISS: wirkt eine Kraft? Inertialsystem?
(begründet), (c) Actio = Reactio in eigenen Worten, plus
Freikörperbild Pferd-Wagen.

\subsection*{Rechnung}
\textbf{a) Die drei Axiome.}
\begin{itemize}
  \item \textbf{1. (Trägheit):} Ein Körper behält (ohne resultierende
        Kraft) seinen Bewegungszustand bei -- Ruhe oder gleichförmige
        Bewegung.
  \item \textbf{2. (Grundgleichung):} Eine resultierende Kraft
        beschleunigt den Körper: $\vec F=m\vec a$, die Beschleunigung
        zeigt in Richtung der Kraft.
  \item \textbf{3. (Actio = Reactio):} Wirkt ein Körper auf einen zweiten
        eine Kraft aus, so wirkt der zweite eine gleich grosse,
        entgegengesetzte Kraft auf den ersten zurück.
\end{itemize}

\textbf{b) Astronaut in der ISS.} Auf den Astronauten wirkt durchaus eine
Kraft: die \emph{Gravitationskraft} der Erde ($\vec F_G$, nach unten/zentral
zur Erde). Er fällt zusammen mit der Station beständig um die Erde. Die
ISS ist \textbf{kein Inertialsystem}: Sie laeuft auf einer gekrümmten
Bahn, ist also beschleunigt ($\vec a\ne0$, Richtung Erde). ``Schwerelos''
heisst hier nur, dass auf den Astronauten keine \emph{Normalkraft} (etwa
vom Boden) wirkt -- die Gewichtskraft ist unverändert da.

\textbf{c) Actio = Reactio und Pferd-Wagen.} Beim Pferd, das einen Wagen
zieht, wirkt auf den Wagen die Kraft $\vec F_{\text{Pferd}\to\text{Wagen}}$
nach vorne; auf das Pferd wirkt zurück die gleich grosse Kraft
$\vec F_{\text{Wagen}\to\text{Pferd}}=-\vec F_{\text{Pferd}\to\text{Wagen}}$
nach hinten. Diese beiden Kräfte heben sich \emph{nicht} gegeneinander auf,
weil sie auf \emph{verschiedene} Körper wirken. Das Pferd kann den Wagen
trotzdem beschleunigen, weil auf das Pferd (ausser der Gegenkraft des
Wagens) die Kraft des Bodens (Reibung der Hufe auf dem Boden) nach vorne
wirkt; am Wagen setzt sich die Zugkraft gegen die Trägheit durch.

\begin{ergebnis}
(a) Trägheit, Grundgleichung $\vec F=m\vec a$, Kraft und Gegenkraft;
(b) Gewichtskraft wirkt, ISS ist kein Inertialsystem, Schwerelosigkeit =
fehlende Normalkraft; (c) Actio/Reactio trifft verschiedene Körper,
Pferd wird durch Bodenreibung nach vorne beschleunigt.
\end{ergebnis}

\begin{kontrolle}
\textbf{Logik:} In (b) wird unterschieden zwischen ``Kraft wirkt'' und
``Normalkraft fehlt''; in (c) wird der Fehler ``Kräfte heben sich auf''
identifiziert. Beide stimmen mit 2. und 3. Axiom überein. \checkmark
\end{kontrolle}

\begin{merksatz}
Kraftwirkung und Gegenkraft wirken auf verschiedene Körper; ein Körper
beschleunigt wegen der resultierenden \emph{Summe} aller auf \emph{ihn}
wirkenden Kräfte. Schwerelosigkeit ist fehlende Normalkraft, nicht
fehlende Gravitation.
\end{merksatz}
